\section{CM twist transport and the exact missing host interface}
\label{sec:cm-twist-transport}

The fixed Hom realization of AC-G23 admits a geometric transport
statement.  It preserves the geometric CM data, with an explicit
finite-prime cost.  This is a theorem about polarized varieties and
coefficient lattices; its hypotheses and output are distinct from
the archived predicate \texttt{HostExW}.

\begin{proposition}[AC-G24: polarized CM transport with one new tame prime]
\label{AC-G24}
Let $F$ be a number field containing $K=\mathbb Q(\omega)$.
Let $(A,\lambda_A)/F$ be a polarized abelian variety of dimension
$g$, with a full integral $\mathcal O_C$ action for a CM field $C$,
defined over $F$.
Suppose $C$ contains $K$ and the Rosati involution of $\lambda_A$
restricts to complex conjugation on $C$.
For a character $\xi:G_F\to\mu_3$, viewed as scalar CM
automorphisms, there is a polarized twist
$(A^\xi,\lambda_A^\xi)/F$ with the following properties.
\begin{enumerate}
\item Over the field cut out by $\xi$, the polarized varieties
are isomorphic.  They have the same dimension, polarization type,
geometric CM type and stable Faltings height.  The full CM action
is still defined over $F$.
\item For every coefficient prime $\ell$ there is an integral
identity of Galois representations with CM action
\[
 T_\ell A^\xi\simeq
 T_\ell A\otimes_{\mathcal O_K\otimes\mathbb Z_\ell}
       (\mathcal O_K\otimes\mathbb Z_\ell)(\xi),        \tag{G24.1}
\]
The integral $\mathcal O_K$ action in this identity is part of
the stated $\mathcal O_C$ action.
\item Put $\xi=\chi_{269}|_{G_F}$ and $n=2g$.
For $v\nmid269$, the local inertia representations of $A$ and
$A^\xi$ are the same.  At $v\mid269$, using any auxiliary
$\ell\ne269$, their Artin conductors satisfy
\[
 a_v(V_\ell A^\xi)\leq a_v(V_\ell A)+n.                \tag{G24.2}
\]
If $V_\ell A$ is unramified and $\xi$ is nontrivial on inertia,
the new exponent is exactly $n$, and the new local Euler factor
is $1$.  In particular, writing $D=[F:\mathbb Q]$,
\[
 \log\operatorname{Norm}_{F/\mathbb Q}\mathfrak f(A^\xi)
 \leq\log\operatorname{Norm}_{F/\mathbb Q}\mathfrak f(A)
             +nD\log269.                               \tag{G24.3}
\]
The corresponding normalized logarithmic increase is at most
$n\log269$.
\end{enumerate}

If a rigidified ample line bundle $\mathcal L/F$ has an
isomorphism $[\omega]^*\mathcal L\simeq\mathcal L$ over $F$
respecting its rigidification at zero, it also descends to a line
bundle $\mathcal L^\xi$ on $A^\xi$.  Its spaces of sections are
obtained by the induced semilinear twist.  Any specified
$\mu_3$-invariant Hermitian metric on these sections is preserved
over the splitting field.  This last assertion is conditional on
the stated line-bundle and metric data; it does not assert an
integral theta-lattice comparison at the ramified prime.
\end{proposition}
\begin{proof}
Let $L/F$ be a finite Galois extension killing $\xi$.
On $A_L$ replace its ordinary semilinear descent action by
\[
             \sigma*x=[\xi(\sigma)]\,\sigma(x).
\]
The automorphisms $[\xi(\sigma)]$ are defined over $F$,
commute, and satisfy the cocycle law; hence this is a descent datum.
Here effectivity can be seen directly using projectivity.  Start
with an ample line bundle on $A_L$ and tensor all its translates
under this finite semilinear group.  The result has a canonical
linearization by permutation of the tensor factors.  A sufficiently
large tensor power gives an equivariant projective embedding.
Finite Galois descent for its finite-dimensional section space and
for every graded piece of its homogeneous ideal then constructs a
projective variety over $F$ with the required base change.  The
identity, multiplication and inverse commute with the descent datum
and descend as morphisms.  Smoothness and connectedness are checked
after $L/F$, so the descended variety is an abelian variety.

For $u\in\mu_3\subset C$, Rosati conjugation gives
$u^\dagger u=\bar u u=1$.  Equivalently
$u^\vee\lambda_Au=\lambda_A$.  Thus the polarization morphism
descends with exactly the same type.  Every element of $C$ commutes
with the cocycle, so its endomorphism descends too.  The geometric
CM type is unchanged under an isomorphism preserving the action.
The equality of stable Faltings heights follows from that polarized
isomorphism over a finite extension.  Applying the covariant Tate
module functor to the displayed semilinear action proves (G24.1).
This uses an actual integral automorphism, and is valid at $\ell=3$;
it requires no averaging projector.

For the local statements, AC-K04 and AC-G22 show that
$\chi_{269}$ is unramified away from $269$ and is actually
trivial at the prime above $3$.  These properties persist under
restriction to $F$.  At a place above $269$ the character is tame,
possibly trivial after local base change.  It is trivial on every
positive upper ramification group, so tensoring leaves the Swan
conductor unchanged.  The tame term is a codimension of inertia
invariants, between zero and $n$; it can therefore increase by at
most $n$.  This proves (G24.2).  If the original representation is
unramified and the new character is nontrivial, a generator of its
inertia image acts through the two CM embeddings of $K$ by
$\omega$ and $\omega^2$.  Neither eigenvalue is one.  Thus the
invariant space is zero, the conductor is $n$, and the determinant
on the invariant space is the constant polynomial $1$.

If $f_v$ and $e_v$ are the residue degree and ramification index
of $v\mid269$, then
$\sum_{v\mid269}e_vf_v=D$ and $\sum_{v\mid269}f_v\leq D$.
Summing the local increase times
$\log\operatorname{Norm}v=f_v\log269$ proves (G24.3).

For the line bundle, choose the stated isomorphism preserving the
rigidification.  Its composite around the third power of
$[\omega]$ is an automorphism of $\mathcal L$.  Since $A$ is
proper and geometrically connected, such an automorphism is a
scalar; its value on the rigidified fiber is one.  The composite
is therefore the identity, giving a genuine $\mu_3$-linearization.
Combining it with the Galois cocycle descends $\mathcal L$ and
its sections.  An invariant metric is unchanged by every element
of this linearization, proving the asserted metric statement.
\end{proof}

\begin{corollary}[A fixed twofold enlargement when the old CM field lacks $\mu_3$]
\label{AC-G24-double}
Let $(A,\lambda_A)/F$ be principally polarized of dimension $g$
with its CM action defined over $F$, and assume $F\supset K$.
Assume its polarization has Rosati involution equal to conjugation
on the given CM field $C$.
The product $B=A^2$ has an integral $\mathcal O_K$ action generated
by
\[
 U=\begin{pmatrix}0&-1\\1&-1\end{pmatrix}.
\]
The polarization obtained from the product polarization by
\[
 H=\begin{pmatrix}2&-1\\-1&2\end{pmatrix}               \tag{G24.4}
\]
is preserved by $U$.  Its type is
$(1,\ldots,1,3,\ldots,3)$ with $g$ entries of each kind,
and its degree is $3^{2g}$.  Twisting this polarized $B$ by
$\chi_{269}$ is therefore possible without putting the splitting
field $K(\sqrt[3]{269})$ inside a CM field.
If the original CM field $C$ does not contain $K$, the commuting
actions give full CM by the field $CK$ on $B$ and on its twist.
The new dimension is $2g$ and the stable Faltings height is
$2h_{\rm Fal}^{\rm st}(A)$.
\end{corollary}
\begin{proof}
Direct multiplication gives $U^2+U+1=0$ and
$U^{\mathsf t}HU=H$.  The positive definite matrix $H$ has
eigenvalues $1,3$.  The associated symmetric positive endomorphism
of the product defines the stated polarization.  Its elementary
divisors are $1,3$, since its entries have gcd one and determinant
three.  On the product integral symplectic lattice these elementary
divisors each occur $g$ times, giving the asserted type and degree.
The matrix identity proves preservation by $U$, so the preceding
descent proof applies.

The scalar action of $C$ on both factors commutes with $U$.
If $K\not\subset C$, then $C\cap K=\mathbb Q$ and
$C\otimes_{\mathbb Q}K=CK$ is a CM field of degree
$2[C:\mathbb Q]=4g$, the required full degree for $A^2$.
The tensor product embedding into $M_2(C)$ is injective, and all
its elements commute with the cocycle and hence descend.  The
Rosati involution is conjugation on each commuting factor; it
remains so after twisting.  Finally stable Faltings height is
additive on products and invariant under geometric isomorphism.
\end{proof}

\begin{remark}[The principal-lattice qualification]
The twofold construction has controlled dimension and polarization
cost, but its polarization is not principal at $3$.  Replacing it
by a principal polarization would require an appropriate isotropic
subgroup of its polarization kernel with descent data.  Neither a
rational quotient nor a Fourier projector supplies that subgroup
for free over the desired field.  Thus (G24.4) is not a replacement
for the principal source lattice $AA=\mathfrak b^{-1}\mathfrak c_E$.
The original AM-H17 CM field
$\mathbb Q(\zeta_{2^{e+2}},\zeta_7,\sqrt{-q},\sqrt{-d})$
does not contain $K$ when $q,d\ne3$: it is unramified at $3$,
whereas $K$ is ramified.  The corollary gives an actual geometric
alternative after a base extension by $K$ of degree at most two,
with the displayed polarization cost, rather than asserting that
its original scalar $\mu_3$ action was already present.
\end{remark}

\begin{proposition}[AC-G25: the precise budget and tying requirements]
\label{AC-G25}
The archived structure \texttt{Rows.Data} alone gives no bound on
its \texttt{shaLen} from the Hermitian spaces, theta maps, lattices,
or idempotents that occur elsewhere in that structure.  More
precisely, for any instance $D$ and any function
$a:\mathbb N^2\to\mathbb Q_{\geq0}$, there is an instance $D_a$
with all the other fields identical and
$D_a.\mathrm{shaLen}(\ell,e)=a(\ell,e)$.

For a fixed $\epsilon>0$, an arithmetic coefficient-transport
estimate in the \emph{same additive units as the source budget}
would be sufficient if it had the form
\[
 W_\xi(x)\leq W_0(x)+C_{269}(1+x)^b,\qquad
 W_0(x)\leq\tfrac12\exp(x^\epsilon/2),                \tag{G25.1}
\]
with $b\geq0$ and $C_{269}>0$ independent of the varying host
and $x$.
Then $W_\xi(x)\leq\exp(x^\epsilon/2)$ for all sufficiently
large $x$ on the same supplied windows.
Neither the finite-conductor estimate (G24.3) nor the logarithmic
coordinate-height estimate (G23.3) proves (G25.1).
\end{proposition}
\begin{proof}
In the exact supplied file
\texttt{Formalism/LocalRingCardinality/DataRows.lean}, the last
two fields of \texttt{Rows.Data} are a function
\texttt{shaLen} with codomain $\mathbb Q$ and its pointwise
nonnegativity proof.  No preceding field mentions that function.
Replace just those two fields by $a$ and its nonnegativity proof.
This is the required $D_a$; all geometric and metric fields remain
identical.  If $t>1$, take $a(3,1)$ to be any desired nonnegative
integer and all its other values to be zero.  The definition
\texttt{Data.VV} then gives $D_a.\mathrm{VV}(3)=a(3,1)$.
Thus even its first prime contribution is arbitrarily large with
those other fields fixed.  This is a countermodel to an inference
from this structure alone, not to a later theorem whose hypotheses
explicitly tie \texttt{shaLen} to arithmetic cohomology.

For (G25.1),
$\log(2C_{269})+b\log(1+x)-x^\epsilon/2\to-\infty$.
Hence $C_{269}(1+x)^b\leq\tfrac12\exp(x^\epsilon/2)$
for all large $x$, and addition proves the assertion.
The normalization qualification is essential.  For example,
multiplication by $269^d$ on a one-dimensional Hermitian space
has scalar of absolute logarithmic height
$h([269^d:1])=d\log269$ and operator norm $269^d$.
If $d=\lceil x^b\rceil$ with $b>\epsilon$, that norm exceeds
$\exp(x^\epsilon/2)$ for all sufficiently large $x$.
Thus a polynomial logarithmic-height bound does not by itself
bound the unlogged norm by the source's allowance.
\end{proof}

\paragraph{Exact archived interface, checked in the supplied source.}
In \texttt{DataHost.lean}, \texttt{primeClass} is defined by the
old ramified classes followed by the Jacobi symbol and
\texttt{H.psAt}; \texttt{FactClassFrob} specifies their Frobenius
coordinates.  The displayed \texttt{AA} is built from the chosen
different-splitting ideal $\mathfrak b$ and
\texttt{H.ccE}.  A coefficient ramified at $269$ introduces a new
local condition even when that prime is unramified in the old
host field.  At a prime where the old representation is unramified
and the new character is not, (G24.2) gives an explicit obstruction
to retaining the old unramified local factor.  A new ramified-prime
clause must use the full inertia-invariant local factor.

The exact file
\path{Development/AutomorphicArithmetic/CharacterBoxes/ExistsPageN.lean}
defines \texttt{VVle E t B} for a general number field, but
\texttt{VAt} and \texttt{VWinAt} supply it only through
\texttt{IsChiField}.  That predicate requires an embedding of
$E$ in a rational cyclotomic field with a kernel described by a
Dirichlet character and two Jacobi symbols.
The same dependency is preserved by \texttt{supplyOfAdm} in
\texttt{HostOfField.lean} and \texttt{supplyChi} in
\texttt{SupplyChi.lean}.  None of these declarations
replaces the growing field $H.E$ by the fixed field $K$ merely
because a coefficient has been realized by CM curves.

The file
\path{Development/AbelianVarietyTheory/Definitions/AutomorphicMembership/IsTwIffCanonEq.lean}
defines \texttt{IsTw} by two identities for the already constructed
\texttt{cst} and \texttt{idx}, and proves
\texttt{isTw\_iff\_canon\_eq} by resetting those two fields.
It is a canonical local-character normalization inside the given
\texttt{TC} data.  It is not an arbitrary Artin-coefficient tensor
transport theorem.

Consequently a narrower geometric route is now concrete: use the
integral CM coefficient of AC-G23, or the polarized twist of AC-G24,
and prove its compatibility with the actual theta lattices and
local conditions.  The required arithmetic estimate must still
hold over the actual varying fields, or come with a proved
reduction of those fields' twisted groups to the fixed exponents
$m_e$ with a controlled error in the units of (G25.1).
A polynomial estimate in the weight over a single fixed field
would be useful input; it does not satisfy that missing
field-uniform assertion by itself.
